\section{Causal-Role Erasure}
\label{sec:problem-formulation}

Object-centered erasure typically asks whether a named object remains visible.
For a physical event, this endpoint is incomplete: the object's visibility may
weaken while its effects remain legible later in the video. We therefore
define the intervention at the level of an entity's role in a registered event
and distinguish the evidence to suppress from the content that must be
preserved.

\subsection{Source, Receiver, and Causal Footprint}
\label{subsec:event-anatomy}

\paragraph{Registered event anatomy.}
We represent case $i$ by the structured binding
\begin{equation}
    e_i=(m_i,s_i,a_i,r_i,f_i,h_i,c_i),
    \label{eq:event-binding}
\end{equation}
where $m_i$ is a mechanism family, $s_i$ is the \emph{source}, $a_i$ is
the registered trigger or interaction, $r_i$ is the \emph{receiver}, $f_i$
is the downstream \emph{causal footprint}, $h_i$ is the registered no-source
receiver state, and $c_i$ contains the remaining scene context. The source is
the entity whose participation initiates the
registered interaction; the receiver is the entity or medium on which that
interaction acts. The footprint is not arbitrary co-occurring content. It is
the mechanism-specific visual consequence expected after the trigger and
dependent on the source's participation---for example, a displaced receiver,
new fragments, or a persistent indentation. The state \(h_i\), such as
remaining upright, intact, or unloaded, specifies what the receiver should do
without the registered event. A structured factual renderer
\(T_{m_i}^{\mathrm f}\) expresses the event binding as prompt \(p_i^{f}\);
its target-side counterpart uses \(h_i\) to construct the no-source supervision in
\S\ref{subsec:counterfactual-pairs}.

\paragraph{Examples and scope.}
In \emph{water impact}, a compact object is the source, its visible entry
through the water surface is the trigger, the water receiver is the
receiver, and the resulting splash and expanding ripples form the footprint.
In \emph{brittle fracture}, an impactor is the source, its downward impact is
the trigger, the initially intact brittle object is the receiver, and the new
broken fragments form the footprint. Our frozen scope further includes rigid
collision, powder impact, elastic deformation, material release, and surface
trace; complete operational definitions appear in
\appref{app:task-ontology}. These examples also separate identity from role:
the same noun may denote a causal source in one prompt and a noncausal
bystander or unrelated object in another. Although a receiver is a distinct
role in the general event definition, the frozen specificity benchmark tests
the protected source noun as a bystander rather than as the receiver.

\subsection{The Role-Conditioned Source-Free World}
\label{subsec:no-source-world}

Erasing a noun unconditionally cannot satisfy this role distinction. It may
under-erase a causal prompt by leaving the source-dependent footprint, or
over-erase a noncausal prompt by suppressing the same object in a legitimate
bystander or otherwise noncausal use. The desired intervention must instead
depend on the source slot occupied by the entity.

For a factual binding $e_i$, the corresponding source-free target is the
conditional video distribution in which the causal source, its registered
trigger, and its dependent footprint are absent, while the receiver, context,
and plausible temporal variation remain. A sample from this
\emph{distributional source-free target} need not match any factual
realization frame by frame, nor recover the unique physical history that would
have occurred without the source. In particular, the source-free video should
remain a natural dynamic video rather than a repeated clean frame. We denote
the prompt specifying this target by $p_i^{cf}$ and construct it in
\S\ref{subsec:counterfactual-pairs}.

This target induces two complementary behavioral contracts. For a causal
prompt $p_i^{f}$, the edited generator should move toward the source-free
distribution by reducing evidence of both $s_i$ and $f_i$ while preserving
$r_i$ and a usable video. For a specificity prompt, the same noun may be
requested in a registered noncausal role, or a visually related footprint may
be legitimately produced by an acceptable alternative cause. The edited
generator should then preserve the protected object and follow that noncausal
request. Thus the task construction and evaluation are defined relative to an
externally registered event binding, not to a global blacklist of nouns or
visual effects. At inference, the generator receives only the rendered text
prompt and the mechanism-specific adapter; it receives no separate event
tuple or role label.

The role labels and mechanism definitions are provided by the structured
dataset. We test whether editing induces role-conditioned behavior; we do not
claim that the model discovers causal roles or learns a true internal causal
representation.

\subsection{Three Desiderata}
\label{subsec:desiderata}

A valid causal-role erasure operator must jointly satisfy three desiderata.
\emph{Efficacy} requires less visible evidence of the registered source and
its footprint on causal cases. \emph{Specificity} requires retaining the same
nouns and legitimate alternative footprints on noncausal cases.
\emph{Validity} requires preserving the receiver and avoiding unusable or
low-quality videos. The last condition prevents scene destruction from being
counted as erasure.

We encode these requirements with full-video atomic scores on a frozen
$0/1/2$ rubric. For causal case $u$ and method $k$, let $S_k(u)$ and $F_k(u)$
denote source and footprint visibility, where $0$, $1$, and $2$ mean absent,
partial, and clear; let $R_k(u)$ and $Q_k(u)$ denote receiver preservation and
video quality, where higher is better. A case earns causal-erasure credit only
when the corresponding unedited backbone $b$ first demonstrates the event:
\begin{align}
E_b(u)={}&\mathbb{1}\!\left[S_{\operatorname{Orig}(b)}(u)=2,
F_{\operatorname{Orig}(b)}(u)\geq 1,
R_{\operatorname{Orig}(b)}(u)\geq 1,
Q_{\operatorname{Orig}(b)}(u)\geq 1\right], \\
U_k(u)={}&\mathbb{1}\!\left[R_k(u)\geq 1,\;Q_k(u)\geq 1\right].
\label{eq:eligibility-usability}
\end{align}
The per-case Causal Erasure Score is then
\begin{equation}
\ces_k(u)=E_b(u)U_k(u)
\frac{(2-S_k(u))+(2-F_k(u))}{4}.
\label{eq:ces}
\end{equation}
Original-ineligible and unusable cases receive zero rather than disappearing
from the denominator. Consequently, \ces{} cannot reward a backbone that never
generated the registered event or an edit that destroys the video.
It is nevertheless an absolute gated output score rather than an improvement
score. Treatment efficacy is established by a paired \ces{} contrast against
the same-backbone Matched Control, or by a registered Original-normalized
contrast for an external method; nonzero credit for one output does not by
itself demonstrate an editing gain.

For a specificity case, let $P_k(u)$ be protected-object visibility and
$N_k(u)$ be adherence to the registered noncausal role, both scored from $0$
to $2$. With
$U_k^{\mathrm{sp}}(u)=\mathbb{1}[R_k(u)\geq1,\,Q_k(u)\geq1]$, Specificity
Utility is
\begin{equation}
\su_k(u)=U_k^{\mathrm{sp}}(u)\frac{P_k(u)+N_k(u)}{4}.
\label{eq:su}
\end{equation}
Both metrics lie in $[0,1]$, and neither replaces the separately reported
receiver, quality, usability, and base-capability statistics. Exact rubrics,
aggregation, uncertainty estimation, and human canonicalization are specified
in \appref{app:evaluation-metrics}.

These scores deliberately admit partial progress. Moving a source or
footprint from clear to partial visibility earns intermediate credit when the
output remains usable; complete absence is not required for nonzero per-case
credit. Source absence, footprint absence, and strict success---which requires
$E_b=1$, $S_k=F_k=0$, and $R_k=Q_k=2$---are therefore secondary outcomes
rather than the definition of success.
